\documentclass{beamer}

\usepackage{amsmath,amsfonts,amssymb,bm}
\usepackage{quantikz}
\usepackage{graphicx}
\usepackage{xcolor}
\usepackage{tikz}

\usetheme{Madrid}

\title{HHL as Spectral Filtering}
\subtitle{Quantum Phase Estimation, Green's Functions, and TDHF Connections}
\author{Morten Hjorth-Jensen}
\date{Spring 2026}

\begin{document}

\frame{\titlepage}

%================================================
\section{Motivation}
%================================================

\begin{frame}{Linear Systems in Physics}
Many problems in physics reduce to solving
\[
A x = b.
\]

Examples:
\begin{itemize}
\item discretized differential equations,
\item response theory,
\item Green's function methods,
\item time-dependent Hartree--Fock (TDHF),
\item random phase approximation (RPA).
\end{itemize}

The HHL algorithm aims to prepare a quantum state proportional to
\[
|x\rangle \propto A^{-1}|b\rangle.
\]
\end{frame}

%------------------------------------------------

\begin{frame}{The Core Question}
Suppose
\[
A|u_j\rangle = \lambda_j |u_j\rangle,
\qquad
|b\rangle = \sum_j \beta_j |u_j\rangle.
\]
Then
\[
A^{-1}|b\rangle
=
\sum_j \frac{\beta_j}{\lambda_j}|u_j\rangle.
\]

Thus HHL is fundamentally about implementing the spectral map
\[
\lambda_j \mapsto \frac{1}{\lambda_j}.
\]

This is why HHL is best viewed as a \emph{spectral filtering algorithm}.
\end{frame}

%================================================
\section{Spectral Decomposition}
%================================================

\begin{frame}{Spectral Representation of the Solution}
Let
\[
A = \sum_j \lambda_j |u_j\rangle \langle u_j|.
\]
Then
\[
A^{-1} = \sum_j \frac{1}{\lambda_j} |u_j\rangle \langle u_j|.
\]
Applying this to the source state gives
\[
|x\rangle \propto A^{-1}|b\rangle
=
\sum_j \frac{\beta_j}{\lambda_j}|u_j\rangle.
\]

The whole algorithm reduces to:
\begin{enumerate}
\item resolve the eigenvalues,
\item invert them,
\item reconstruct the filtered state.
\end{enumerate}
\end{frame}

%------------------------------------------------

\begin{frame}{Why Quantum Phase Estimation Appears}
Quantum phase estimation (QPE) is needed in the original HHL algorithm because it allows us to attach eigenvalue information to the eigenvectors:
\[
|u_j\rangle |0\rangle
\longrightarrow
|u_j\rangle |\lambda_j\rangle.
\]

After QPE, the state becomes
\[
\sum_j \beta_j |u_j\rangle |\lambda_j\rangle,
\]
so that a controlled operation can implement the map
\[
\lambda_j \mapsto \frac{1}{\lambda_j}.
\]
\end{frame}

%================================================
\section{The HHL Circuit}
%================================================

\begin{frame}{Overall Structure of HHL}
The original HHL algorithm has four main ingredients:
\begin{enumerate}
\item state preparation,
\item quantum phase estimation,
\item controlled eigenvalue inversion,
\item inverse phase estimation.
\end{enumerate}

The key spectral information enters through the QPE block.
\end{frame}

%------------------------------------------------

\begin{frame}{HHL Circuit: High-Level View}
\[
\begin{quantikz}[row sep=0.35cm, column sep=0.45cm]
\lstick{$|b\rangle$} & \qw & \qw & \gate[wires=2,style={fill=blue!15,draw=blue}]{\text{QPE}} & \qw & \qw & \gate[wires=2,style={fill=yellow!20,draw=orange}]{\text{Controlled } \lambda^{-1}} & \qw & \gate[wires=2,style={fill=blue!15,draw=blue}]{\text{QPE}^\dagger} & \qw & \qw \\
\lstick{$|0\rangle^{\otimes m}$} & \qw & \qw & & \qw & \qw & & \qw & & \qw & \qw \\
\lstick{$|0\rangle_a$} & \qw & \qw & \qw & \qw & \qw & \gate{R_y(\lambda^{-1})} & \qw & \qw & \meter{} & \qw
\end{quantikz}
\]

\begin{itemize}
\item Blue block: \textbf{QPE} extracts eigenvalues.
\item Yellow block: controlled rotation implements inversion.
\item Final measurement post-selects the desired branch.
\end{itemize}
\end{frame}

%------------------------------------------------

\begin{frame}{Where the QPE Sits}
The essential block is
\[
\begin{quantikz}[row sep=0.4cm, column sep=0.5cm]
\lstick{$|u_j\rangle$} & \gate[wires=2,style={fill=blue!15,draw=blue}]{\text{QPE}} & \qw \\
\lstick{$|0\rangle^{\otimes m}$} & & \rstick{$|\lambda_j\rangle$}
\end{quantikz}
\]
which performs
\[
|u_j\rangle |0\rangle^{\otimes m}
\longrightarrow
|u_j\rangle |\lambda_j\rangle.
\]

Without this step, the algorithm does not know which spectral component to invert.
\end{frame}

%------------------------------------------------

\begin{frame}{Detailed HHL Circuit}
\[
\begin{quantikz}[row sep=0.35cm, column sep=0.38cm]
\lstick{$|b\rangle$} 
& \qw 
& \gate[wires=2,style={fill=blue!15,draw=blue}]{\text{QPE}} 
& \qw 
& \qw 
& \gate[wires=2,style={fill=yellow!20,draw=orange}]{\text{Controlled Rotation}} 
& \qw
& \gate[wires=2,style={fill=blue!15,draw=blue}]{\text{QPE}^\dagger} 
& \qw 
& \qw \\
\lstick{$|0\rangle^{\otimes m}$} 
& \qw 
& 
& \qw 
& \qw 
& 
& \qw
& 
& \qw 
& \qw \\
\lstick{$|0\rangle_a$} 
& \qw 
& \qw 
& \qw 
& \qw 
& \gate{R_y(\theta(\lambda))} 
& \qw
& \qw 
& \meter{} 
& \qw
\end{quantikz}
\]

Here the angle is chosen so that
\[
\sin\frac{\theta(\lambda_j)}{2} \propto \frac{C}{\lambda_j}.
\]
\end{frame}

%================================================
\section{Step-by-Step State Evolution}
%================================================

\begin{frame}{Step 1: Prepare the Input State}
Start with
\[
|b\rangle = \sum_j \beta_j |u_j\rangle.
\]

The ancilla and clock/eigenvalue registers are initialized as
\[
|b\rangle |0\rangle^{\otimes m} |0\rangle_a.
\]
\end{frame}

%------------------------------------------------

\begin{frame}{Step 2: Apply QPE}
After QPE:
\[
|b\rangle |0\rangle^{\otimes m}|0\rangle_a
\longrightarrow
\sum_j \beta_j |u_j\rangle |\lambda_j\rangle |0\rangle_a.
\]

This is the crucial spectral decomposition stage.
\end{frame}

%------------------------------------------------

\begin{frame}{Step 3: Controlled Rotation}
Apply a controlled rotation on the ancilla:
\[
|\lambda_j\rangle |0\rangle_a
\longrightarrow
|\lambda_j\rangle
\left(
\sqrt{1-\frac{C^2}{\lambda_j^2}}\,|0\rangle_a
+
\frac{C}{\lambda_j}|1\rangle_a
\right).
\]

Hence the global state becomes
\[
\sum_j \beta_j |u_j\rangle |\lambda_j\rangle
\left(
\sqrt{1-\frac{C^2}{\lambda_j^2}}\,|0\rangle_a
+
\frac{C}{\lambda_j}|1\rangle_a
\right).
\]
\end{frame}

%------------------------------------------------

\begin{frame}{Step 4: Uncompute the Eigenvalue Register}
Apply inverse QPE:
\[
\sum_j \beta_j |u_j\rangle |\lambda_j\rangle
\left(
\sqrt{1-\frac{C^2}{\lambda_j^2}}\,|0\rangle_a
+
\frac{C}{\lambda_j}|1\rangle_a
\right)
\]
\[
\longrightarrow
\sum_j \beta_j |u_j\rangle
\left(
\sqrt{1-\frac{C^2}{\lambda_j^2}}\,|0\rangle_a
+
\frac{C}{\lambda_j}|1\rangle_a
\right).
\]

Now the spectral information has been transferred into the ancilla amplitudes.
\end{frame}

%------------------------------------------------

\begin{frame}{Step 5: Postselection}
Conditioning on measuring the ancilla in the state \(|1\rangle_a\), the system register collapses to
\[
|x\rangle \propto \sum_j \frac{\beta_j}{\lambda_j}|u_j\rangle
= A^{-1}|b\rangle.
\]

Thus HHL realizes the inverse operator spectrally.
\end{frame}

%================================================
\section{HHL as Spectral Filtering}
%================================================

\begin{frame}{HHL as an Operator Filter}
The central transformation is
\[
A = \sum_j \lambda_j |u_j\rangle\langle u_j|
\quad \longrightarrow \quad
A^{-1} = \sum_j \frac{1}{\lambda_j}|u_j\rangle\langle u_j|.
\]

This is a spectral filter:
\[
f(\lambda) = \frac{1}{\lambda}.
\]

Therefore HHL is not merely a solver of linear equations. It is a quantum method for applying a function of an operator to a state.
\end{frame}

%------------------------------------------------

\begin{frame}{Comparison with Time Evolution}
Compare:
\[
e^{-iAt} = \sum_j e^{-i\lambda_j t}|u_j\rangle\langle u_j|
\]
with
\[
A^{-1} = \sum_j \frac{1}{\lambda_j}|u_j\rangle\langle u_j|.
\]

Time evolution multiplies each eigenmode by a phase.
HHL multiplies each eigenmode by an inverse eigenvalue.

So HHL is closer in spirit to:
\begin{itemize}
\item Green's functions,
\item resolvents,
\item response theory,
\item inverse-operator methods.
\end{itemize}
\end{frame}

%================================================
\section{Green's Functions and Resolvents}
%================================================

\begin{frame}{Resolvent Formulation}
In many-body physics, one often studies the resolvent
\[
G(z) = (zI - H)^{-1}.
\]

If
\[
H|u_j\rangle = E_j |u_j\rangle,
\]
then
\[
G(z) = \sum_j \frac{1}{z-E_j}|u_j\rangle\langle u_j|.
\]

This has exactly the same spectral structure as HHL.
\end{frame}

%------------------------------------------------

\begin{frame}{Applying a Green's Function to a Source State}
Let
\[
|b\rangle = \sum_j \beta_j |u_j\rangle.
\]
Then
\[
G(z)|b\rangle
=
\sum_j \frac{\beta_j}{z-E_j}|u_j\rangle.
\]

This is identical in form to the HHL output
\[
A^{-1}|b\rangle
=
\sum_j \frac{\beta_j}{\lambda_j}|u_j\rangle.
\]

Thus HHL can be viewed as a Green's-function-style spectral inverse algorithm.
\end{frame}

%------------------------------------------------

\begin{frame}{Retarded Green's Function}
In response theory one often uses
\[
G^R(\omega) = \frac{1}{\omega + i\eta - H}.
\]

Spectrally,
\[
G^R(\omega)
=
\sum_j
\frac{1}{\omega + i\eta - E_j}
|u_j\rangle\langle u_j|.
\]

This shows that the inverse operator acts as a frequency-dependent filter:
\begin{itemize}
\item states near resonance are enhanced,
\item off-resonant states are suppressed.
\end{itemize}
\end{frame}

%------------------------------------------------

\begin{frame}{Dynamical Susceptibility}
A typical response function has the form
\[
\chi(\omega)
=
\langle \psi_0 |
\hat{O}^\dagger
\frac{1}{\omega + i\eta - (H-E_0)}
\hat{O}
|\psi_0\rangle.
\]

Defining
\[
|b\rangle = \hat{O}|\psi_0\rangle,
\]
we obtain
\[
\chi(\omega)
=
\langle b|
\frac{1}{\omega + i\eta - (H-E_0)}
|b\rangle.
\]

Hence HHL is naturally related to the computation of response states and susceptibilities.
\end{frame}

%================================================
\section{Connection to TDHF and Linear Response}
%================================================

\begin{frame}{TDHF Equation}
Time-dependent Hartree--Fock is governed by
\[
i \frac{d\rho}{dt} = [h[\rho],\rho].
\]

Linearizing around a stationary state \(\rho_0\),
\[
\rho(t) = \rho_0 + \delta\rho(t),
\]
one obtains an equation for the fluctuation \(\delta \rho\).
\end{frame}

%------------------------------------------------

\begin{frame}{Linearized TDHF in Frequency Space}
After linearization and Fourier transform, one gets a linear-response equation of the form
\[
(\omega I - \mathcal{L}) \delta \rho(\omega) = s(\omega),
\]
where
\begin{itemize}
\item \(\mathcal{L}\) is the linearized Liouvillian,
\item \(s(\omega)\) is the external driving term.
\end{itemize}

This is exactly a linear system:
\[
A x = b
\]
with
\[
A = \omega I - \mathcal{L}.
\]
\end{frame}

%------------------------------------------------

\begin{frame}{TDHF as an HHL-Type Problem}
The formal correspondence is
\[
A x = b
\qquad \Longleftrightarrow \qquad
(\omega I - \mathcal{L}) \delta \rho(\omega) = s(\omega).
\]

Thus
\[
\delta \rho(\omega)
=
(\omega I - \mathcal{L})^{-1} s(\omega).
\]

This is precisely the same inverse-operator structure implemented by HHL.
\end{frame}

%------------------------------------------------

\begin{frame}{Connection to RPA}
In the small-amplitude limit, TDHF is closely related to the random phase approximation (RPA), often written as
\[
\begin{pmatrix}
A & B \\
-B^* & -A^*
\end{pmatrix}
\begin{pmatrix}
X \\ Y
\end{pmatrix}
=
\omega
\begin{pmatrix}
X \\ Y
\end{pmatrix}.
\]

Equivalently, one may formulate a driven response equation
\[
(\omega I - M) v = s.
\]

Again, the physical problem becomes spectral inversion.
\end{frame}

%================================================
\section{Why QPE is Essential in Original HHL}
%================================================

\begin{frame}{Why QPE is Needed}
In the original HHL algorithm, QPE is the mechanism that resolves the operator spectrally:
\[
|u_j\rangle \mapsto |u_j\rangle |\lambda_j\rangle.
\]

Only then can the ancilla rotation implement
\[
\frac{1}{\lambda_j}.
\]

So the logical structure is:
\begin{enumerate}
\item diagonalize in the eigenbasis,
\item apply the spectral filter,
\item reconstruct the filtered state.
\end{enumerate}
\end{frame}

%------------------------------------------------

\begin{frame}{Physical Interpretation of QPE in HHL}
From a physics viewpoint:
\begin{itemize}
\item QPE plays the role of spectral analysis,
\item controlled rotation applies the resolvent-like weight,
\item inverse QPE recombines the filtered modes.
\end{itemize}

Therefore:
\[
\text{HHL} = \text{spectral decomposition} + \text{inverse weighting} + \text{recombination}.
\]

This is why HHL is naturally understood through Green's functions and linear response.
\end{frame}

%================================================
\section{Caveats and Remarks}
%================================================

\begin{frame}{Important Caveats}
The connection is mathematically deep, but practical implementation faces challenges:
\begin{itemize}
\item efficient Hamiltonian simulation or block encoding,
\item conditioning of the matrix,
\item non-Hermitian response matrices in some TDHF/RPA settings,
\item state preparation and measurement costs.
\end{itemize}

Still, the conceptual bridge is very strong:
\[
\text{linear systems} \leftrightarrow \text{resolvents} \leftrightarrow \text{response theory}.
\]
\end{frame}

%------------------------------------------------

\begin{frame}{Hermitian Embedding for Non-Hermitian Problems}
If the response matrix is non-Hermitian, one may embed it into a Hermitian operator:
\[
\mathcal{H}
=
\begin{pmatrix}
0 & M \\
M^\dagger & 0
\end{pmatrix}.
\]

This allows one to reformulate a non-Hermitian inversion problem in a larger Hermitian space, which is often the natural entry point for HHL-like methods.
\end{frame}

%================================================
\section{Summary}
%================================================

\begin{frame}{Summary}
\begin{itemize}
\item In original HHL, \textbf{QPE is essential}.
\item It is the block that extracts the eigenvalues \(\lambda_j\).
\item HHL implements the spectral filter
\[
\lambda_j \mapsto \frac{1}{\lambda_j}.
\]
\item This makes HHL naturally connected to
\begin{itemize}
\item Green's functions,
\item resolvents,
\item TDHF,
\item RPA,
\item linear response theory.
\end{itemize}
\end{itemize}
\end{frame}

%------------------------------------------------

\begin{frame}{Final Conceptual Picture}
\[
\boxed{
\text{HHL}
=
\text{Quantum Phase Estimation}
+
\text{Spectral Inversion}
+
\text{State Reconstruction}
}
\]

\[
\boxed{
A^{-1}|b\rangle
\sim
(zI-H)^{-1}|b\rangle
\sim
(\omega I-\mathcal{L})^{-1} s
}
\]

So HHL is best understood as a quantum inverse-operator method.
\end{frame}

\end{document}
